
\subsubsection{Research Questions}

\textbf{RQ1 (Existence):} Do duality equations produce valid, non-divergent SSM parameters?

\textbf{RQ2 (Mechanism):} Does duality-based initialization provide lower reconstruction error than random?

\subsubsection{Dataset}

WikiText-103 \citep{merity2017wikitext} validation split: 100 samples for stability testing (h-e1), 500 samples for reconstruction comparison (h-m1), with 64--512 tokens per sample.

\subsubsection{Models}

\textbf{Teacher:} BERT-base-uncased (12 layers, 768 hidden dimension, 12 attention heads, approximately 110M parameters)

\textbf{Student:} SSM with d\_state = 64, d\_model = 768

\subsubsection{Evaluation Metrics}

\begin{table}[htbp]
\centering
\resizebox{\columnwidth}{!}{%
\begin{tabular}{lll}
\toprule
Hypothesis & Metric & Threshold \\
\midrule
h-e1 & NaN/Inf Rate & 0\% \\
h-e1 & Magnitude Ratio & < $10\times$ \\
h-m1 & Reconstruction Error & Duality < Random \\
h-m1 & P-value & < 0.05 \\
\bottomrule
\end{tabular}
}
\end{table}

\subsubsection{Statistical Analysis}

Paired t-tests compare duality and random initialization errors across the 500 samples. Cohen's d quantifies effect size. Per-layer analysis is conducted across all 12 BERT layers.
